%% TC adaptation: runtime stake -> feedback-write failure -> mechanism -> evidence.

Accelerator-backed generative services usually treat a returned sequence as an
endpoint. Revision-generative ASR does not: Whisper-style decoders condition on
previously generated text~\cite{radford2023whisper,machacek2023whisperstreaming},
and Qwen3-ASR re-prompts itself with rollback prefixes from earlier
revisions~\cite{shi2026qwen3asr}. A generated hypothesis therefore mutates the
state of later execution. Yet throughput-oriented runtimes optimize request
scheduling and output delivery without exposing an authorization boundary
between model generation and this feedback write~\cite{kwon2023vllm,yu2022orca,zheng2024sglang}.

Batching improves accelerator utilization by sharing execution across
sessions~\cite{kwon2023vllm,agrawal2024sarathi,zhong2024distserve,sun2024llumnix},
but it can change floating-point reduction order and hence a greedy token
sequence relative to singleton execution~\cite{yuan2025numerical}. A terminal
service contains that difference within one response. In a revision-generative
service, the differing output can become the parent of every later revision
(Fig.~\ref{fig:motivation}). Controlled paired interventions establish that
this is a feedback-channel failure, not merely display variation: 8 of 21
legal state pairs propagate when the feedback edge is present, none propagate
when it is removed, and the measured revision amplification factor (RAF)
reaches 153.

Existing mechanisms regulate different boundaries. LocalAgreement stabilizes
display prefixes but does not authorize feedback parents~\cite{machacek2023whisperstreaming};
speculative decoding verifies draft tokens against a target
distribution~\cite{leviathan2023specdec,chen2023specsample}; and LLM-42 provides
verify--rollback for deterministic token generation, not revision-state
installation~\cite{zhang2026llm42}. Numerical reproducibility repairs execution
rather than state provenance~\cite{yuan2025numerical}. Optimistic concurrency
control supplies a useful analogy~\cite{kung1981occ,bernstein1983mvcc}, but here
the speculative write is a model-generated projection whose legality must be
established by the serving runtime.

We accordingly treat each revision output as a speculative state write:
$\mathrm{Display}(\out_t)\nRightarrow
\mathrm{Authorize}(\proj(\out_t;H_t))$. We introduce \textbf{Dual-Granularity
Feedback Transactions} (DGFT), which decouples a fine-grained proposal/display
path from a coarse canonical authority. Only a certificate bound to the parent
version, input frontier, projection, and execution contract can install state;
same-granularity certification is the $\granratio=1$ case. \textbf{RevTrace}
implements DGFT as a two-plane parallel runtime with contract-scoped
certification. The abstraction targets feedback-driven generation, while all
empirical generality claimed here is explicitly limited to the evaluated
revision-generative ASR pipelines.

This paper contributes:
\begin{enumerate}
  \item a formal feedback-state contamination model and causal evidence that
    identifies the feedback edge as the propagation channel
    (Sections~\ref{sec:model} and~\ref{sec:eval});
  \item DGFT and the RevTrace runtime, with certificate-, version-, and
    frontier-bound commits, three state-history invariants, and a cost model
    for selecting authority granularity
    (Sections~\ref{sec:dgft}--\ref{sec:impl}); and
  \item a contract-bounded evaluation on Qwen3-ASR 0.6B/1.7B and Whisper
    tiny/base/small.en, using two RTX~3090 GPUs and an independent
    8$\times$V100 platform. It covers matched-workload throughput, display and
    commit latency, energy, family-specific quality, protocol fuzzing, and the
    measured single-node capacity boundary (Section~\ref{sec:eval}).
\end{enumerate}
